% \iffalse meta-comment
%
% hit-thesis-geometry.dtx 是版心尺寸
%
% \fi
%
% \section{版心尺寸}
%
% ^^A ======================================================================
% ^^A Module thesis-geometry: 版芯尺寸与页边距（hit-thesis）
% ^^A ======================================================================
%    \begin{macrocode}
%<*thesis-geometry>
%    \end{macrocode}
%
% \changes{v3.2a}{2026/08/14}{版心分成新旧两套：新版面按 Word 页面设置算，是主路；
%   旧版面收进一个可整块删除的段落，由 \option{v2-layout} 要回来}
% 版心有两套。\emph{新版面}按现行规范排，取值从学校发的 Word 范例里解出来，
% 走 \cs{hit@msword@layout:n} 换算，这是主路，什么都不写就是它。\emph{旧版面}是
% v3.1e 及更早的那几组硬编码尺寸，靠类选项 \option{v2-layout} 显式要回来。
%
% 下面「旧版面」那一整段是补丁：停止维护旧分支时，连同
% \file{src/setup/hit-thesis-options.dtx} 里的 \option{v2-layout} 声明一起删掉即可，
% 新版面那一段不依赖它。
%
% 两套共用的只有纸张与四个全局开关。
%    \begin{macrocode}
\bool_if:NTF \g__hit_debug_bool
  { \RequirePackage [ showframe ] { geometry } }
  { \RequirePackage { geometry } }
\geometry { a4paper , hcentering , ignoreall , nomarginpar }
%    \end{macrocode}
%
% \subsection{新版面}
%
% \changes{v3.2a}{2026/08/14}{本科接上新版面：取值走 \option{layout} 与
%   \option{format} 两族，写在 \file{src/setup/hit-defaults.dtx} 里}
% 新版面一个数都不写在这里。取值是 \option{layout} 与 \option{format} 两族的
% 键值表，与用户能写的完全同一套，摆在 \file{src/setup/hit-defaults.dtx} 的
% \cs{hitpresetup} 队列里；换算规则见 \file{src/utils/hit-msword.dtx} 与
% \file{src/utils/hit-format.dtx}。
%
% 队列的回放晚于本文件（\pkg{docstrip} 的顺序见 \file{src/hit-thesis.dtx} 的装配表），
% 所以新版面这一路在这里什么也不做，只是把旧版面那一段让开。\pkg{geometry} 的
% 版心是在 \cs{begin}\marg{document} 时才算的，回放的位置来得及。
%
% 本科与硕博走这条路，共用一套取值。硕博的图题表题间距与本科有出入，那一处
% 还没接，接的时候只加 \option{caption-figure} 与 \option{caption-table} 两个目标，
% 别的不动。博士后还留在旧版面：那一档另有一套参数，还没量。
% 威海本科没有对照文档，跟着哈深那一份走，没有逐页核对过。
%
% \subsection{旧版面（补丁，可整块删除）}
%
% 退回这里：本科写 \option{v2-layout}\texttt{=\{bachelor=true\}}，硕博写
% \option{v2-layout}\texttt{=\{graduate=one\}} 或 \texttt{two}；博士后眼下
% 只有这一条路。
%    \begin{macrocode}
\bool_if:NT \g__hit_layout_two_bool
  {
%    \end{macrocode}
%
% \changes{v2.0.0}{2018/06/14}{添加版芯设置选项}
% \changes{v3.2a}{2026/08/14}{删掉 \option{newgeometry}\texttt{!=no} 那一支版心。
%   它是 PlutoThesis 时代留下的，与现行规范差得太远，选项取值一并删除}
% 两组尺寸，由 \option{v2-layout} 的 \option{graduate} 选：\texttt{two} 是版心
% 236\,mm 高，\texttt{one} 是 240\,mm。默认走 \texttt{two}。
%    \begin{macrocode}
    \bool_if:NTF \g__hit_layout_two_graduate_two_bool
      {
        \geometry
          {
            centering ,
            text = { 150true~mm , 236true~mm } ,
            left = 30true~mm , head = 5true~mm , headsep = 2true~mm ,
            footskip = 0true~mm , foot = 5.2true~mm
          }
      }
      {
        \geometry
          {
            centering ,
            text = { 150true~mm , 240true~mm } ,
            left = 30true~mm , head = 5true~mm , headsep = 0true~mm ,
            footskip = 0true~mm , foot = 0true~mm
          }
      }
%    \end{macrocode}
%
% \changes{v3.1d}{2025/03/03}{为哈尔滨本科特意改了页眉位置。\issue[172]}
% \changes{v3.2a}{2026/08/16}{这一组去掉校区限定，本科三校区已统一}
% 本科的旧版面另有一组页眉位置与版心高度。这一条属于旧版面那一套，
% 新版面里本科的页眉位置由 Word 设置里的「页眉距边界」算出来，不需要特例。
%
% \changes{v3.2a}{2026/08/08}{哈尔滨本科的版心高度改写成 \option{textheight}，
%   原先写 \option{bottom}，与上面的 \option{text} 一起构成三个纵向约束，
%   \pkg{geometry} 只能丢掉一个并报 Over-specification}
% 纵向只能给两个约束（另一个是纸高）。上面 \option{text} 已经给了高度，这里再给
% \option{top} 与 \option{bottom} 就是三个，\pkg{geometry} 会丢掉 \option{text}
% 那个高度、按 \texttt{297-36.5-28.8} 算出 231.7\,mm，并报一条警告。直接写
% \option{textheight} 说的是同一件事，警告没了。
%
% 高度写成 \texttt{659.25035pt} 而不是 \texttt{231.7true\,mm}：两者只差末位，
% 但 \pkg{geometry} 的换算路径不同，写毫米会让一个 PDF 跳转锚点的纵坐标差
% 0.01bp。写 \pkg{geometry} 自己算出来的那个值，生成物逐字节不变。
%    \begin{macrocode}
    \bool_if:NT \g__hit_bachelor_bool
      {
        \geometry
          { top = 36.5true~mm , headsep = 1true~mm , textheight = 659.25035pt }
      }
  }
%</thesis-geometry>
%    \end{macrocode}
%
% \Finale
